\begin{table}[htbp]
\centering
\caption{Tổng hợp một số hướng nghiên cứu liên quan}
\label{tab:ch1-nghien-cuu}
\footnotesize
\begin{tabularx}{\textwidth}{@{}l l Y@{}}
\toprule
\textbf{Công trình} & \textbf{Hệ thống và phương pháp} & \textbf{Điểm liên quan và giới hạn so với luận văn} \\
\midrule
Mao và cộng sự \cite{Mao2018RSMABridging, Mao2022RSMAFundamentals}
& RSMA nhiều người dùng; phân tích và tối ưu
& Cung cấp nền tảng luồng chung và luồng riêng; không tập trung STAR--RIS và biểu diễn hành động DRL. \\
\addlinespace
Xu, Mu và cộng sự \cite{Xu2021STARRIS, Mu2021STARAided}
& STAR--RIS; mô hình hệ thống ES, MS, TS
& Cung cấp mô hình truyền và phản xạ; không giải bài toán DRL của luận văn. \\
\addlinespace
Li và cộng sự \cite{Li2022RSMARISInterplay}
& RIS kết hợp RSMA; phân tích tương tác
& Làm rõ tiềm năng kết hợp hai công nghệ; không tập trung cấu trúc SISO và bộ giải mã học. \\
\addlinespace
Meng và cộng sự \cite{Meng2024PPOSTARRSMA}
& STAR--RIS hỗ trợ RSMA; PPO
& Cho thấy PPO có thể dùng cho bài toán liên tục; luận văn mở rộng đánh giá sang DDPG, TD3 và nhóm AO trên cùng mô hình. \\
\addlinespace
Zhong và cộng sự \cite{Zhong2022HybridRLSTAR}
& STAR--RIS; học tăng cường lai
& Khai thác cấu trúc điều khiển STAR--RIS; mục tiêu và mô hình khác RSMA SISO của luận văn. \\
\addlinespace
Ge và cộng sự \cite{Ge2024RSMAMassiveMIMO}
& STAR--RIS, massive MIMO, RSMA; tối ưu chung
& Hệ nhiều anten và mục tiêu khác; không phản ánh đặc điểm không có tạo búp sóng của SISO. \\
\bottomrule
\end{tabularx}
\tablesource{Tác giả tổng hợp từ các công trình trích dẫn (2026)}
\end{table}
